% Options for packages loaded elsewhere
\PassOptionsToPackage{unicode}{hyperref}
\PassOptionsToPackage{hyphens}{url}
\PassOptionsToPackage{dvipsnames,svgnames,x11names}{xcolor}
\documentclass[
  10pt,
  fontset=fandol]{ctexart}
\usepackage{xcolor}
\usepackage[margin=16mm]{geometry}
\usepackage{amsmath,amssymb}
\setcounter{secnumdepth}{-\maxdimen} % remove section numbering
\usepackage{iftex}
\ifPDFTeX
  \usepackage[T1]{fontenc}
  \usepackage[utf8]{inputenc}
  \usepackage{textcomp} % provide euro and other symbols
\else % if luatex or xetex
  \usepackage{unicode-math} % this also loads fontspec
  \defaultfontfeatures{Scale=MatchLowercase}
  \defaultfontfeatures[\rmfamily]{Ligatures=TeX,Scale=1}
\fi
\usepackage{lmodern}
\ifPDFTeX\else
  % xetex/luatex font selection
\fi
% Use upquote if available, for straight quotes in verbatim environments
\IfFileExists{upquote.sty}{\usepackage{upquote}}{}
\IfFileExists{microtype.sty}{% use microtype if available
  \usepackage[]{microtype}
  \UseMicrotypeSet[protrusion]{basicmath} % disable protrusion for tt fonts
}{}
\makeatletter
\@ifundefined{KOMAClassName}{% if non-KOMA class
  \IfFileExists{parskip.sty}{%
    \usepackage{parskip}
  }{% else
    \setlength{\parindent}{0pt}
    \setlength{\parskip}{6pt plus 2pt minus 1pt}}
}{% if KOMA class
  \KOMAoptions{parskip=half}}
\makeatother
\setlength{\emergencystretch}{3em} % prevent overfull lines
\providecommand{\tightlist}{%
  \setlength{\itemsep}{0pt}\setlength{\parskip}{0pt}}
\usepackage{amsmath,amssymb,mathtools,needspace}
\xeCJKDeclareCharClass{CJK}{"2032,"0400->"04FF,"2160->"216B,"2460->"2473,"25A0,"25C7,"25CB,"25CF}
\IfFontExistsTF{Arial Unicode MS}{\xeCJKsetup{AutoFallBack=true}\setCJKfallbackfamilyfont{\CJKrmdefault}{Arial Unicode MS}}{}
\setcounter{secnumdepth}{0}
\setlength{\emergencystretch}{2em}
\allowdisplaybreaks
\usepackage{bookmark}
\IfFileExists{xurl.sty}{\usepackage{xurl}}{} % add URL line breaks if available
\urlstyle{same}
\hypersetup{
  pdftitle={用正交相似变换约化矩阵},
  colorlinks=true,
  linkcolor={Maroon},
  filecolor={Maroon},
  citecolor={Blue},
  urlcolor={Blue},
  pdfcreator={LaTeX via pandoc}}

\title{用正交相似变换约化矩阵}
\author{}
\date{}

\begin{document}
\maketitle

\subsubsection{9.3.2 用正交相似变换约化矩阵}\label{ux7528ux6b63ux4ea4ux76f8ux4f3cux53d8ux6362ux7ea6ux5316ux77e9ux9635}

下面考虑用初等反射阵来正交相似约化一般矩阵和对称矩阵。设

\href{../../source-images/pdf-246.jpeg}{原页图像}

\[
A=\left(\begin{array}{c|ccc}
a_{11}&a_{12}&\cdots&a_{1n}\\\hline
a_{21}&a_{22}&\cdots&a_{2n}\\
\vdots&\vdots&&\vdots\\
a_{n1}&a_{n2}&\cdots&a_{nn}
\end{array}\right)
\equiv\begin{pmatrix}
a_{11}&A_{12}^{(1)}\\
a_{21}^{(1)}&A_{22}^{(1)}
\end{pmatrix},
\]

\textbf{步 1} 不妨设 \(a_{21}^{(1)}\ne0\)，否则这一步不需约化，选择初等反射阵 \(R_1\) 使 \(R_1a_{21}^{(1)}=-\sigma_1e_1\)，其中

\[
\left\{
\begin{aligned}
\sigma_1&=\operatorname{sgn}(a_{21})\left(\sum_{i=2}^na_{i1}^2\right)^{\frac12},\\
u_1&=a_{21}^{(1)}+\sigma_1e_1,\\
\rho_1&=\frac12\|u_1\|_2^2=\sigma_1(\sigma_1+a_{21}),\\
R_1&=I-\rho_1^{-1}u_1u_1^{\mathrm T}.
\end{aligned}
\right.
\tag{9.3.1}
\]

令 \(U_1=\begin{pmatrix}I&0\\0&R_1\end{pmatrix}\)，则

\[
A_2=U_1A_1U_1=\begin{pmatrix}
a_{11}&A_{21}^{(1)}R_1\\
R_1a_{21}^{(1)}&R_1A_{22}^{(1)}R_1
\end{pmatrix}
\equiv\begin{pmatrix}
A_{11}^{(2)}&a_{12}^{(2)}&A_{13}^{(2)}\\
O&a_{22}^{(2)}&A_{23}^{(2)}
\end{pmatrix},
\]

其中 \(A_{11}^{(2)}\in\mathbf R^{2\times1},\allowbreak \quad a_{22}^{(2)}\in\mathbf R^{n-2},\allowbreak \quad A_{23}^{(2)}\in\mathbf R^{(n-2)\times(n-2)}\)。

\textbf{步 \(k\)} 设对 \(A\) 已进行了第 \(k-1\) 步正交相似约化，即 \(A_k\) 有形式

\[
\begin{aligned}
A_k&=U_{k-1}A_{k-1}U_{k-1}\\
&=\left(\begin{array}{ccc|c|ccc}
a_{11}&a_{12}^{(2)}&\cdots&a_{1k}^{(k)}&a_{1,k+1}^{(k)}&\cdots&a_{1n}^{(k)}\\
-\sigma_1&a_{22}^{(2)}&\cdots&a_{2k}^{(k)}&a_{2,k+1}^{(k)}&\cdots&a_{2n}^{(k)}\\
&\ddots&\ddots&\vdots&\vdots&&\vdots\\
&&-\sigma_{k-1}&a_{kk}^{(k)}&a_{k,k+1}^{(k)}&\cdots&a_{k,n}^{(k)}\\\hline
&&&a_{k+1,k}^{(k)}&a_{k+1,k+1}^{(k)}&\cdots&a_{k+1,n}^{(k)}\\
&&&\vdots&\vdots&&\vdots\\
&&&a_{nk}^{(k)}&a_{n,k+1}^{(k)}&\cdots&a_{nn}^{(k)}
\end{array}\right)\\
&\equiv\begin{pmatrix}
A_{11}^{(k)}&a_{12}^{(k)}&A_{13}^{(k)}\\
O&a_{22}^{(k)}&A_{23}^{(k)}
\end{pmatrix},
\end{aligned}
\]

其中 \(A_{11}^{(k)}\in\mathbf R^{k\times(k-1)},\allowbreak \quad a_{22}^{(k)}\in\mathbf R^{n-k},\allowbreak \quad A_{23}^{(k)}\in\mathbf R^{(n-k)\times(n-k)}\)。

设 \(a_{22}^{(k)}\ne0\)，选择初等反射阵 \(R_k\)，使 \(R_ka_{22}^{(k)}=-\sigma_ke_1\)，其中

\[
\left\{
\begin{aligned}
\sigma_k&=\operatorname{sgn}(a_{k+1,k}^{(k)})\left(\sum_{i=k+1}^na_{ik}^2\right)^{\frac12},\\
u_k&=a_{22}^{(k)}+\sigma_ke_1,\\
\rho_k&=\frac12\|u_k\|_2^2=\sigma_k(\sigma_k+a_{k+1,n}^{(k)}),\\
R_k&=I-\rho_k^{-1}u_ku_k^{\mathrm T}.
\end{aligned}
\right.
\tag{9.3.2}
\]

\begin{quote}
\textbf{校注（agent 补充，ch09a-E009）} 本页 \(A_2\) 的二阶分块表达式右上块原印 \(A_{21}^{(1)}R_1\)，已照录。由页首定义的右上块 \(A_{12}^{(1)}\) 和 \(U_1\) 的分块乘法，此处应为 \(A_{12}^{(1)}R_1\)。见\href{../assets/ch09a/review-p246-initial-partition.png}{页首分块}和\href{../assets/ch09a/review-p246-A-2.png}{放大原式}。

\textbf{校注（agent 补充，ch09a-E010）} 式（9.3.2）原图 \(\sigma_k\) 求和内为 \(a_{ik}^2\)，未印迭代上标；\(\rho_k\) 最后一个矩阵元为 \(a_{k+1,n}^{(k)}\)，均照录。由本页 \(a_{22}^{(k)}=(a_{k+1,k}^{(k)},\ldots,a_{nk}^{(k)})^{\mathrm T}\) 和 \(u_k=a_{22}^{(k)}+\sigma_ke_1\)，求和应使用当前列元素 \((a_{ik}^{(k)})^2\)，而范数平方恒等式给出 \(\rho_k=\sigma_k(\sigma_k+a_{k+1,k}^{(k)})\)。故前者疑漏迭代上标，后者列下标疑将 \(k\) 印成 \(n\)。见\href{../assets/ch09a/review-p246-9-3-2.png}{放大原式}。
\end{quote}

\href{../../source-images/pdf-247.jpeg}{原页}

设 \(U_k=\begin{pmatrix}I&O\\O&R_k\end{pmatrix}\)，则

\[
A_{k+1}=U_kA_kU_k=
\begin{pmatrix}
A_{11}^{(k)}&a_{12}^{(k)}&A_{13}^{(k)}R_k\\
O&R_ka_{22}^{(k)}&R_kA_{23}^{(k)}R_k
\end{pmatrix}
=
\begin{pmatrix}
A_{11}^{(k)}&a_{12}^{(k)}&A_{13}^{(k)}R_k\\
O&-\sigma_ke_1&R_kA_{23}^{(k)}R_k
\end{pmatrix}.
\tag{9.3.3}
\]

由式（9.3.3）知，\(A_{k+1}\) 的左上角 \(k+1\) 阶子阵为上 Hessenberg 阵，从而约化又进了一步，重复这过程，直到

\[
U_{n-2}\cdots U_2U_1AU_1U_2\cdots U_{n-2}
=\begin{pmatrix}
a_{11}&\times&\times&\cdots&\times\\
-\sigma_1&a_{22}^{(2)}&\times&\cdots&\times\\
&-\sigma_2&a_{33}^{(3)}&\ddots&\vdots\\
&&\ddots&\ddots&\times\\
&&&-\sigma_{n-1}&a_{nn}^{(n-1)}
\end{pmatrix}=A_{n-1}.
\]

总结上述讨论，有如下定理。

\textbf{定理 9.13}　如果 \(A\in\mathbb R^{n\times n}\)，则存在初等反射阵 \(U_1,U_2,\cdots,U_{n-2}\)，使

\[
U_{n-2}\cdots U_2U_1AU_1U_2\cdots U_{n-2}=C\quad\text{（上 Hessenberg 阵）}.
\]

在 \(A_k\to A_{k+1}=U_kA_kU_k\) 的进一步约化中，需要计算 \(R_k\) 和 \(A_{13}^{(k)}R_k\)，\(R_kA_{23}^{(k)}R_k\)。用初等反射阵正交相似约化 \(A\) 为上 Hessenberg 阵，大约需要 \(\dfrac53n^3\) 次乘法运算。

由于 \(U_k\) 都是正交阵，所以 \(A_1\sim A_2\sim\cdots\sim A_{n-1}\)。求 \(A\) 的特征值问题，就转化为求上 Hessenberg 阵 \(C\) 的特征值问题。由定理 9.13，记 \(P=U_{n-2}\cdots U_2U_1\)，则

\[
PAP^{\mathrm T}=C.
\]

设 \(y\) 是 \(C\) 的对应特征值 \(\lambda\) 的特征向量，则 \(P^{\mathrm T}y\) 为 \(A\) 的对应特征值 \(\lambda\) 的特征向量，且

\[
P^{\mathrm T}y=U_1U_2\cdots U_{n-2}y
=(I-\lambda_1^{-1}u_1u_1^{\mathrm T})\cdots(I-\lambda_{n-2}^{-1}u_{n-2}u_{n-2}^{\mathrm T})y.
\]

\textbf{定理 9.14}　如果 \(A\in\mathbb R^{n\times n}\) 为对称阵，则存在初等反射阵 \(U_1,U_2,\cdots,U_{n-2}\)，使

\[
U_{n-2}\cdots U_2U_1AU_1U_2\cdots U_{n-2}=A_{n-1}
=\begin{pmatrix}
c_1&b_1&&&\\
b_1&c_2&b_2&&\\
&\ddots&\ddots&\ddots&\\
&&b_{n-2}&c_{n-1}&b_{n-1}\\
&&&b_{n-1}&c_n
\end{pmatrix}\equiv C.
\]

\textbf{证明}　由定理 9.13，存在初等反射阵 \(U_1,U_2,\cdots,U_{n-2}\)，使 \(A_{n-1}\) 为上 Hessenberg 阵，但 \(A_{n-1}\) 又为对称阵，因此 \(A_{n-1}\) 为对称三对角阵。

由上面讨论可知，在由 \(A_k\to A_{k+1}=U_kA_kU_k\) 一步计算过程中，只需计算 \(R_k\) 和 \(R_kA_{23}^{(k)}R_k\)。由于 \(A\) 的对称性，故只需计算 \(R_kA_{23}^{(k)}R_k\) 的对角线下面的元素。注意到

\[
R_kA_{23}^{(k)}R_k=(I-\rho_k^{-1}u_ku_k^{\mathrm T})(A_{23}^{(k)}-\rho_k^{-1}A_{23}^{(k)}u_ku_k^{\mathrm T}),
\]

\begin{quote}
校注（agent 补充）：本页 \(P^{\mathrm T}y\) 展开式原印为 \(\lambda_1^{-1}\)、\(\lambda_{n-2}^{-1}\)，而同页下方及次页初等反射阵的归一化因子使用 \(\rho_k^{-1}\)。此处保留原印 \(\lambda\)，记为疑似原书符号不一致。
\end{quote}

\href{../../source-images/pdf-248.jpeg}{原页}

引进记号

\[
r_k=\rho_k^{-1}A_{23}^{(k)}u_k,\qquad
 t_k=r_k-\frac{\rho_k^{-1}}2(u_k^{\mathrm T}r_k)u_k,
\]

则

\[
R_kA_{23}^{(k)}R_k=A_{23}^{(k)}-u_kt_k^{\mathrm T}-t_ku_k^{\mathrm T}
\quad(i=k+1,\cdots,n;\ j=k+1,\cdots,i).
\]

\textbf{算法 3}　正交相似约化对称阵为对称三对角阵。设 \(A\in\mathbb R^{n\times n}\) 是对称阵，本算法确定初等反射阵 \(U_1,U_2,\cdots,U_{n-2}\)，使 \(U_{n-2}\cdots U_1AU_1\cdots U_{n-2}=C\)（对称三对角阵），\(C\) 的对角元 \(c_i\) 存放在单元 \(c_1,c_2,\cdots,c_n\) 中，\(C\) 的非对角元 \(b_i\) 存放在单元 \(b_1,b_2,\cdots,b_{n-1}\) 中。单元 \(b_1,b_2,\cdots,b_n\) 最初可用来存放 \(r_k\) 及 \(t_k\) 的分量，确定 \(U_k\) 的向量 \(u_k\) 的分量 \(u_{k+1,k},\cdots,u_{nk}\)，存放在 \(A\) 的相应位置。\(\rho_k\) 冲掉 \(a_{kk}\)。约化 \(A\) 的结果冲掉 \(A\)，数组 \(A\) 的上部元素不变。如果步 \(k\) 不需要变换，则置 \(\rho_k\) 为零。

对于 \(k=1,2,\cdots,n-2\)，做到 \(L\)。

\textbf{步 1}　\(c_k=a_{kk}\)。

\textbf{步 2}　确定变换：

（1）计算 \(\eta=\max\limits_{k+1\leq i\leq n}|a_{ik}|\)；

（2）如果 \(\eta=0\)，则

\[
\begin{cases}
a_{kk}\to\rho_k=0,\\
b_k\to0,\\
\text{转 }L,\text{否则继续};
\end{cases}
\]

（3）计算 \(a_{ik}\leftarrow u_{ik}=a_{ik}/\eta\quad(i=k+1,\cdots,n)\)；

（4）\(\sigma=\operatorname{sgn}(u_{k+1,k})\sqrt{u_{k+1,k}^2+\cdots+u_{nk}^2}\)；

（5）\(u_{k+1,k}\leftarrow u_{k+1,k}+\sigma\)；

（6）\(a_{kk}\leftarrow\rho_k=\sigma u_{k+1,k}\)；

（7）\(b_k\leftarrow-\sigma\eta\)。

\textbf{步 3}　应用变换：

（1）\(\sigma=0\)；

（2）计算 \(A_{23}^{(k)}u_k\) 及 \(u_k^{\mathrm T}r_k\)，对于 \(i=k+1,\cdots,n\)，做

\[
\begin{cases}
b_i\leftarrow s=\displaystyle\sum_{j=k+1}^{n}a_{ij}u_{jk}+\displaystyle\sum_{j=i+1}^{n}a_{ji}u_{jk},\\
\sigma\leftarrow\sigma+su_{ik};
\end{cases}
\]

（3）计算 \(t_k\)

\[
b_i\leftarrow\rho_k^{-1}\left(b_i-\rho_k^{-1}\frac\sigma2u_{ik}\right)
\quad(i=k+1,\cdots,n);
\]

（4）计算 \(R_kA_{23}^{(k)}R_k\)

对于 \(i=k+1,\allowbreak k+2,\allowbreak \cdots,\allowbreak n;\ j=k+1,\allowbreak \cdots,\allowbreak i\)，做 \(a_{ij}\leftarrow a_{ij}-u_{ik}b_j-b_iu_{jk}\)。

\(L\)：继续循环。

对于 \(k=n-1\)，有

\[
c_{n-1}\leftarrow a_{n-1,n-1},\qquad c_n\leftarrow a_{nn},\qquad b_{n-1}\leftarrow a_{n,n-1}.
\]

\begin{quote}
校注（agent 补充）：算法 3 步 3（2）第一个求和上限原印为 \(n\)。按只存储、更新下三角部分的算法结构，该上限疑应为 \(i\)，否则 \(j>i\) 的项与第二个求和重叠，且会读入未更新的上三角元素。正文保留原印，列为疑似原书错误。
\end{quote}

\href{../../source-images/pdf-249.jpeg}{原页}

将对称阵 \(A\) 用初等反射阵正交相似约化为对称三对角阵约需做 \(\dfrac23n^3\) 次乘法运算。

用正交矩阵进行约化，有一些特点，如构造的 \(U_k\) 容易求逆，且 \(U_k\) 的元素数量级不大，因此这个算法是十分稳定的。

\textbf{例 9.5}　用 Householder 方法将下述矩阵化为上 Hessenberg 阵

\[
A_1=A=\begin{pmatrix}-4&-3&-7\\2&3&2\\4&2&7\end{pmatrix}.
\]

\textbf{解}　（1）对于 \(k=1\)，确定变换 \(U_1=\begin{pmatrix}1&\begin{matrix}0&0\end{matrix}\\\begin{matrix}0\\0\end{matrix}&R_1\end{pmatrix}\)，\(a_{21}^{(1)}=\begin{pmatrix}2\\4\end{pmatrix}\)，其中 \(R_1\) 为初等反射阵且使

\[
R_1a_{21}^{(1)}=-\sigma_1\begin{pmatrix}1\\0\end{pmatrix},\qquad
\sigma_1=\|a_{21}^{(1)}\|_2=\sqrt{20}\approx4.472\,136,
\]

\[
u_1=a_{21}^{(1)}+\sigma_1\rho_1=\begin{pmatrix}2+\sqrt{20}\\4\end{pmatrix}
\approx\begin{pmatrix}6.472\,136\\4\end{pmatrix},
\]

\[
\rho_1=\sigma_1(\sigma_1+a_{21})=\sqrt{20}(\sqrt{20}+2)\approx28.944\,27,\qquad
R_1=I-\rho_1^{-1}u_1u_1^{\mathrm T}.
\]

（2）计算 \(R_1A_{22}^{(1)}\)。记

\[
A_{22}^{(1)}=\begin{pmatrix}3&2\\2&7\end{pmatrix}\equiv(a_1,a_2),
\]

于是

\[
R_1A_{22}^{(1)}=(R_1a_1,R_1a_2)
=\begin{pmatrix}-3.130\,496&-7.155\,419\\-1.788\,855&1.341\,640\end{pmatrix},
\]

其中

\[
R_1a_i=(I-\rho_1^{-1}u_1u_1^{\mathrm T})a_i
=a_i-(\rho_1^{-1}u_1^{\mathrm T}a_i)u_1\quad(i=1,2).
\]

（3）计算 \(A_{12}^{(1)}R_1\) 及 \((R_1A_{22}^{(1)})R_1\)，即计算

\[
\begin{pmatrix}A_{12}^{(1)}\\(R_1A_{22}^{(1)})\end{pmatrix}R_1
\equiv\begin{pmatrix}b_1^{\mathrm T}\\b_2^{\mathrm T}\\b_3^{\mathrm T}\end{pmatrix}R_1
=\begin{pmatrix}b_1^{\mathrm T}R_1\\b_2^{\mathrm T}R_1\\b_3^{\mathrm T}R_1\end{pmatrix}
=\begin{pmatrix}
7.602\,634&-0.447\,212\\
7.800\,003&-0.399\,999\\
-0.399\,999&2.200\,000
\end{pmatrix},
\]

其中

\[
b_i^{\mathrm T}R_1=b_i^{\mathrm T}-(\rho_1^{-1}b_i^{\mathrm T}u_1)u_1^{\mathrm T}
\quad(i=1,2,3).
\]

（4）计算 \(A_2=U_1A_1U_1\)。

\[
A_2=\begin{pmatrix}-4&A_{12}^{(1)}R_1\\
\begin{matrix}-\sigma_1\\0\end{matrix}&R_1A_{22}^{(1)}R_1
\end{pmatrix}
=\begin{pmatrix}
-4&7.602\,634&-0.447\,212\\
-4.472\,136&7.800\,003&-0.399\,999\\
0&-0.399\,999&2.200\,000
\end{pmatrix}
\]

为上 Hessenberg 阵。

\begin{quote}
校注（agent 补充）：例 9.5（1）原印 \(u_1=a_{21}^{(1)}+\sigma_1\rho_1\) 中的 \(\rho_1\) 已经由放大原图确认。结合紧随的列向量以及同段将 \(\rho_1\) 定义为标量的公式，这里疑应为第一单位向量 \(e_1\)；正文保留原印。
\end{quote}

\subsection{本节覆盖原页的校注记录}\label{ux672cux8282ux8986ux76d6ux539fux9875ux7684ux6821ux6ce8ux8bb0ux5f55}

下列记录按原页关联，含同页相邻小节的校注；计算时同时读取上文校注和完整原页。

\begin{itemize}
\tightlist
\item
  \texttt{ch09a-E009}：原书 PDF 页 246
\item
  \texttt{ch09a-E010}：原书 PDF 页 246
\item
  \texttt{ch09b-E01}：原书 PDF 页 247
\item
  \texttt{ch09b-E02}：原书 PDF 页 248
\item
  \texttt{ch09b-E03}：原书 PDF 页 249
\end{itemize}

\end{document}
